\section{Conclusion}\label{sec:conclusion}
Beyond the technical decisions covered so far, delivering this project also depended on how effectively we worked as a team.
This section reflects on that collaboration, the opportunities we would pursue next, and how well the final result met our original goals.

\subsection{Collaboration}
Building a full-stack project required organized teamwork through \emph{GitHub}.
We used issues to track features, separate branches to work independently, and pull requests to review code before merging.
This helped us catch bugs early and share responsibility across the team.

Because almost every visualization changed multiple times, we focused on keeping our code maintainable.
We wrote modular code, documented key interfaces, and continuously cleaned up our codebase.

Finally, we automated our workflow using \emph{GitHub Actions}.
On every pull request, the pipeline runs backend tests against a temporary \emph{PostgreSQL} database, checks the frontend code, and builds our \LaTeX{} documentation.

\subsection{Future Work}
Throughout the project, we identified several opportunities that would enhance the dashboard's usefulness and analytical depth.

\paragraph{Actual Trip Routes}
Because the dataset records only origin, destination, and duration, exact travel paths are unknown.
GPS traces or routing engines would map trips onto the street network, transforming straight corridors into realistic cycling routes and enriching spatial visualizations.

\paragraph{Station Availability History}
As noted in Section~\ref{sec:backend}, dock occupancy is limited to live snapshots.
Persisting this data over time would reveal availability patterns, showing when shortages occur, how long they last, and how effectively rebalancing responds.

\paragraph{Demand Forecasting}
The current dashboard is entirely descriptive.
However, as our temporal and weather views demonstrate, ridership responds predictably to weather, time, and day.
Adding a forecasting model to these features would predict station-level demand, moving the tool from descriptive analysis to operational planning.

\subsection{Final Considerations}
The project met its main goal.
The interface highlights key patterns immediately, including weekday commutes, rain sensitivity, and differences between members and casual riders.
Throughout the project, we chose clear, practical solutions over complex ones, creating a system that is easy to understand, maintain, and expand.